\chapter*{Extended Abstract}
\addcontentsline{toc}{chapter}{Extended Abstract}

Last-mile delivery is the most expensive and operationally complex segment of the
logistics supply chain, yet small and medium-sized enterprises (SMEs) in Hong Kong
still rely on fragmented, manually coordinated tools for dispatching, navigation,
and proof of delivery.  Commercially available platforms either target large-scale
operators with heavyweight enterprise integrations or offer only basic route
planning without AI assistance, real-time coordination, or multi-tenant security.
This project addresses that gap by designing, implementing, and evaluating
\textbf{OrderVista}, an AI-enhanced logistics dispatch platform that unifies
administrative dispatching and driver-side field operations within a single
cross-platform application.

The system is built with Flutter (targeting Android and web from one codebase),
Supabase (PostgreSQL with Row-Level Security, authentication, real-time
subscriptions, and object storage), Mapbox (forward geocoding and static map
visualisation), and the DeepSeek R1 large language model accessed via the
OpenRouter gateway.  The architecture follows a layered, offline-first design:
a Drift-powered SQLite cache provides instant startup and queued offline
mutations, while a service--repository layer mediates between the local database
and the cloud backend.

Fourteen objectives were defined, spanning six core platform goals and eight
advanced optimisation targets.  All fourteen were achieved:

\begin{itemize}
  \item A \textbf{unified administrative portal} supporting task creation, CSV
        bulk import with geocoding, route generation, driver assignment, proof-of-delivery
        review, fleet management, and real-time chat.
  \item A \textbf{driver mobile application} with optimised task lists, turn-by-turn
        navigation handoff, camera-based proof capture (photo plus barcode/QR scan),
        telemetry reporting, and an AI assistant.
  \item A \textbf{nearest-neighbour route sequencing algorithm} enhanced with
        Haversine distance, mid-shift GPS-based re-sequencing, dynamic priority
        recalculation, and fleet-wide multi-driver distance minimisation, yielding a
        measured 25--31\% reduction in total delivery distance on representative
        test scenarios.
  \item An \textbf{agentic AI assistant} that goes beyond conventional chatbot
        summarisation: the DeepSeek model returns structured JSON actions
        (e.g.\ ``send notification'', ``reassign task'') which are surfaced as
        one-tap execution cards with explicit human-oversight safeguards.  In
        evaluation, the agent correctly identified and structured the appropriate
        action in six out of six test prompts.
  \item \textbf{Enterprise-grade multi-tenant security} via over twenty PostgreSQL
        Row-Level Security policies, a centralised \texttt{is\_company\_member()}
        helper, and a three-role RBAC model (driver, dispatcher, admin).
        Cross-tenant query tests confirmed complete data isolation.
  \item A \textbf{performance-aware driver recommendation engine} combining
        availability, GPS proximity, current workload, completion rate, and average
        delivery time into a composite score computed by a server-side SQL function.
  \item \textbf{Customer-facing delivery tracking} through auto-generated
        tracking codes and a public RPC endpoint requiring no authentication.
\end{itemize}

The system was evaluated through over 26~functional test cases executed across
both admin and driver roles on a physical Android device, supplemented by
unit tests for algorithmic components and dedicated RLS security verification.
Route optimisation was benchmarked against unsequenced baselines; AI quality
was assessed for both summary accuracy and agentic action correctness.
All objectives were mapped to test evidence confirming full compliance.

The principal contributions of this project are threefold.  First, it
demonstrates that agentic AI (where a language model proposes structured
workflow actions under human oversight) is a viable and practical pattern
for logistics decision support, an application area where existing research
has not yet explored LLM-driven action execution.  Second, it shows that
PostgreSQL Row-Level Security can enforce tenant isolation at the database
layer without application-side middleware, offering a lightweight
multi-tenancy model suitable for SME platforms.  Third, it validates an
offline-first mobile architecture with deterministic sync-queue processing
and exponential backoff, ensuring field operability in connectivity-poor
environments common to Hong Kong's dense urban infrastructure.

The signed Android release APK (75.7\,MB) constitutes a fully functional
prototype suitable for demonstration and further development toward
commercial viability.

\newpage
